\begin{example}[Counterfactual query for a random walk]\label{exa:rwCounterfactuals}

    We build a counterfactual twinned network (see \figref{fig:rwCounterfactualNetwork}), by adding counterfactual variables $\tildecatvariableof{[\catorder]}$ linked to the same response variables.
    Further, we include do variables to the counterfactual world, to answer counterfactual queries.

    \begin{figure}[t]
        \stantikzpic{
            \drawvariablenode{6}{2}{$\dovariableof{\catenumerator}$}{DK}
            \drawvariablenode{-3}{0}{$\catvariableof{\catenumerator}$}{XK}
            \drawvariablenode{0}{2}{$\selvariableof{\catenumerator}$}{LK}
            \drawvariablenode{3}{0}{$\tildecatvariableof{\catenumerator}$}{TXK}

            \drawvariablenode{-3}{4}{$\catvariableof{\catenumerator-1}$}{XKmin}
            \drawvariablenode{3}{4}{$\tildecatvariableof{\catenumerator-1}$}{TXKmin}

            \drawtextnode{-2.5}{7}{\small Observed world}
            \draw[dashed] (1.5,8) -- (1.5,-2);
            \drawtextnode{7}{7}{\small Counterfactual world}

            \coordinate (EK) at ($(XK) + (1,2)$);
            \coordinate (TEK) at ($(TXK) + (0,2)$);

            \draw[->-] (XKmin) -- (EK);
            \draw[->-] (LK) -- (EK);
            \draw[->-] (EK) -- (XK);

            \draw[->-] (TXKmin) -- (TEK);
            \draw[->-] (DK) -- (TEK);
            \draw[->-] (LK) -- (TEK);

           \draw[->-] (TEK) -- (TXK);


            \begin{scope}[shift={(10,0)}]

                \draw (0,1) rectangle (6,3);
                \node[anchor=center] at (3,2) {\corelabelsize $\bencodingof{\exfunctionof{\atomenumerator}}$};
                \draw[->-] (3,1) -- (3,-1) node[midway,right] {\colorlabelsize $\catvariableof{\catenumerator}$};
                \draw[->-] (3,5) -- (3,3) node[midway,right] {\colorlabelsize $\catvariableof{\catenumerator-1}$};

                \draw[->-] (8,2) -- (6,2);
                \drawvariabledot{8}{2}
                \draw (6.5,3.5) rectangle (9.5,6.5);
                \drawtextnode{8}{5}{
                    \tiny $\begin{pmatrix}
                               p \\ 1-2p \\ p
                    \end{pmatrix}$}
                \draw [->-] (8,3.5) -- (8,2) node[midway,right] {\colorlabelsize $\selvariableof{\catenumerator}$};
                \draw[->-] (8,2) -- (10,2);


                \begin{scope}[shift={(10,0)}]
                    \draw (0,1) rectangle (6,3);
                    \node[anchor=center] at (3,2) {\corelabelsize $\bencodingof{\causalsymbol,\exfunctionof{\atomenumerator}}$};
                    \draw[->-] (3,1) -- (3,-1) node[midway,right] {\colorlabelsize $\tildecatvariableof{\catenumerator}$};
                    \draw[->-] (3,5) -- (3,3) node[midway,right] {\colorlabelsize $\tildecatvariableof{\catenumerator-1}$};
                    \draw[->-] (8,2) -- (6,2) node[midway,above] {\colorlabelsize $\dovariableof{\catenumerator}$};
                \end{scope}


            \end{scope}
        }
        \caption{Part of the counterfactual twinned network for the random walk (see \exaref{exa:rwResponse} and \exaref{exa:rwCounterfactuals}).}\label{fig:rwCounterfactualNetwork}
    \end{figure}



    Setup:
    \begin{itemize}
        \item The variables $\catvariableof{\catenumerator-1},\catvariableof{\catenumerator},\catvariableof{\catenumerator+1}$ have been observed.
        \item The counterfactual variable $\catvariableof{\catenumerator}$ is intervened on.
    \end{itemize}
    This counterfactual query is answered by the contraction
    \stantikzpic{
        \begin{scope}[shift={(0,4)}]
            \drawtensorblock{3}{2}{3}{$\bencodingof{\exfunctionof{\atomenumerator}}$}
            \drawddwire{3}{3}{$\catvariableof{\catenumerator-1}$}
            %          \drawtextnode{3}{6}{$\vdots$}

            \drawldwire{6}{2}{}
            \drawvariabledot{8}{2}
            \draw (6.5,3.5) rectangle (9.5,6.5);
            \drawtextnode{8}{5}{
                \tiny $\begin{pmatrix}
                           p \\ 1-2p \\ p
                \end{pmatrix}$}
            \draw [->-] (8,3.5) -- (8,2) node[midway,right] {\colorlabelsize $\selvariableof{\catenumerator}$};
            \drawrdwire{10}{2}{}
            \begin{scope}[shift={(10,0)}]
                \drawtensorblock{3}{2}{3}{$\bencodingof{\causalsymbol,\exfunctionof{\atomenumerator}}$}
                \drawddwire{3}{3}{$\tildecatvariableof{\catenumerator-1}$}
                \drawsquaretensor{3}{6}{$\probtensor$}
%                \drawtextnode{3}{6}{$\vdots$}
                \drawldwire{6}{2}{$\dovariableof{\catenumerator}$}
            \end{scope}
        \end{scope}

        \drawtensorblock{3}{2}{3}{$\bencodingof{\exfunctionof{\atomenumerator+1}}$}
        \drawddwire{3}{-1}{$\catvariableof{\catenumerator+1}$}
        %\drawtextnode{3}{-1.5}{$\vdots$}
        \drawddwire{3}{3}{$\catvariableof{\catenumerator}$}

        \drawldwire{6}{2}{}
        \drawvariabledot{8}{2}
        \draw (6.5,0.5) rectangle (9.5,-2.5);
        \drawtextnode{8}{-1}{
            \tiny $\begin{pmatrix}
                       p \\ 1-2p \\ p
            \end{pmatrix}$}
        \draw [->-] (8,0.5) -- (8,2) node[midway,right] {\colorlabelsize $\selvariableof{\catenumerator+1}$};
        \drawrdwire{10}{2}{}

        \begin{scope}[shift={(10,0)}]
            \drawtensorblock{3}{2}{3}{$\bencodingof{\causalsymbol,\exfunctionof{\atomenumerator+1}}$}
            \drawddwire{3}{-1}{$\tildecatvariableof{\catenumerator+1}$}
            %\drawtextnode{3}{-1.5}{$\vdots$}
            \drawddwire{3}{3}{$\tildecatvariableof{\catenumerator}$}
            \drawldwire{6}{2}{$\dovariableof{\catenumerator+1}$}
            \drawsquaretensor{9}{2}{$\onehotmapof{\catdim}$}
        \end{scope}

        \drawsquaretensor{-2}{9}{$\onehotmapof{\catindexof{\catenumerator-1}}$}
        \drawsquaretensor{-2}{4.5}{$\onehotmapof{\catindexof{\catenumerator}}$}
        \drawsquaretensor{-2}{-1}{$\onehotmapof{\catindexof{\catenumerator+1}}$}

        \drawsquaretensor{19}{6}{$\onehotmapof{\tildecatindexof{\catenumerator}}$}

        \draw[->-] (-1,9) -- (3,9);
        \drawvariabledot{3}{9}

        \draw[->-] (-1,4.5) -- (3,4.5);
        \drawvariabledot{3}{4.5}

        \draw[->-] (-1,-1) -- (3,-1);
        \drawvariabledot{3}{-1}

        \drawtextnode{20.5}{4}{{${=}$}}

        \begin{scope}[shift={(17.5,0)}]

            \begin{scope}[shift={(0,4)}]

                \begin{scope}[shift={(10,0)}]
                    \drawtensorblock{3}{2}{3}{$\bencodingof{\causalsymbol,\exfunctionof{\atomenumerator}}$}
                    \drawddwire{3}{3}{$\tildecatvariableof{\catenumerator-1}$}
                    %\drawtextnode{3}{6}{$\vdots$}
                    \drawsquaretensor{3}{6}{$\probtensor$}
                    \drawldwire{6}{2}{$\dovariableof{\catenumerator}$}
                \end{scope}

            \end{scope}

            \begin{scope}[shift={(10,0)}]

                \drawtensorblock{3}{2}{3}{$\bencodingof{\causalsymbol,\exfunctionof{\atomenumerator+1}}$}
                \drawddwire{3}{-1}{$\tildecatvariableof{\catenumerator+1}$}
                %\drawtextnode{3}{-1.5}{$\vdots$}
                \drawddwire{3}{3}{$\tildecatvariableof{\catenumerator}$}
                \drawldwire{6}{2}{$\dovariableof{\catenumerator+1}$}
                \drawsquaretensor{9}{2}{$\onehotmapof{\catdim}$}

            \end{scope}

            \drawsquaretensor{19}{6}{$\onehotmapof{\tildecatindexof{\catenumerator}}$}

            \drawtensorblock{6}{6}{2}{$\onehotmapof{\catindexof{\catenumerator}-\catindexof{\catenumerator-1}+1}$}
            \drawrdwire{10}{6}{$\selvariableof{\catenumerator}$}
            \drawtensorblock{6}{2}{2}{$\onehotmapof{\catindexof{\catenumerator+1}-\catindexof{\catenumerator}+1}$}
            \drawrdwire{10}{2}{$\selvariableof{\catenumerator+1}$}
        \end{scope}

        \drawtextnode{38}{4}{{${=}$}}

        \begin{scope}[shift={(26.5,2)}]
            \drawtensorblock{15}{2}{3}{$\onehotmapof{\tildecatindexof{\catenumerator}+ (\catindexof{\catenumerator+1}-\catindexof{\catenumerator})}$}
            \drawddwire{15}{-1}{$\catvariableof{\catenumerator+1}$};
        \end{scope}
    }
    %If not further interventions are done one the counterfactual world,

    Let us consider an arbitrary set $\arbset\subset[\catorder]$ of variables intervened on.
    For each $\seccatenumerator$ we denote $\catenumerator$ as the previous variable which has been intervened on.

    The counterfactual world is by the above contraction only supported at the indices
    \begin{align*}
        \tildecatindexof{\seccatenumerator}
        = \tildecatindexof{\catenumerator} + \left(\catindexof{\seccatenumerator} - \catindexof{\catenumerator}\right) \, .
    \end{align*}
    If no intervention has been done before $\seccatenumerator$, we define $\tildecatindexof{\catenumerator}=\catindexof{\catenumerator}$ and have $\tildecatindexof{\seccatenumerator}=\catindexof{\seccatenumerator}$.

    % Interpretation
    Thus, the same steps (left, right neighbor or keeping positions) are done in the counterfactual world.
    Given the observed world, the distribution of the counterfactual world variables is therefore deterministic.

    % More general: Non-deterministic counterfactual - Unclear!
%    In more generality, we have a larger support of the response variables after contraction with the one-hot encoded observation.
%    In that case, we would have a non-deterministic behaviour of the counterfactual variables.

%    \begin{align*}
%        \contractionof{
%            \onehotmapofat{\catindexof{\catenumerator}}{\catvariableof{\catenumerator}},
%            \onehotmapofat{\catindexof{\catenumerator-1}}{\catvariableof{\catenumerator-1}},
%            \bencodingofat{}{}
%        }{\selvariableof{\catenumerator}}
%    \end{align*}

\end{example}